\documentclass[aspectratio=169,11pt]{beamer}
\usepackage{paperdeck}
\newcommand{\mgrid}[2]{%
  \begin{tikzpicture}[scale=0.26]
    \foreach \i in {0,...,5}{\foreach \j in {0,...,5}{
      \pgfmathtruncatemacro{\ok}{#2}
      \ifnum\ok=1 \fill[accent] (\j,-\i) rectangle ++(0.9,-0.9);
      \else \fill[black!8] (\j,-\i) rectangle ++(0.9,-0.9); \fi }}
    \node[font=\scriptsize,anchor=north,align=center] at (3,-6.2) {#1};
  \end{tikzpicture}}
\title[T5]{Exploring the Limits of Transfer Learning with a Unified Text-to-Text Transformer}
\subtitle{Raffel et al.\ (JMLR 2020), arXiv:1910.10683}
\author[Paper summary]{Paper summary for the CS-781 capstone}
\date{}
\begin{document}

\begin{frame}\titlepage\end{frame}

\paperglance
{C.\ Raffel, N.\ Shazeer, A.\ Roberts, K.\ Lee, S.\ Narang, M.\ Matena, Y.\ Zhou, W.\ Li, P.\ J.\ Liu. \emph{Exploring the Limits of Transfer Learning with a Unified Text-to-Text Transformer}. JMLR 21, 2020. Google.}
{Cast every NLP task as ``text in, text out'', then run a large controlled study of architectures, objectives, data, fine-tuning and scale. The result is \textbf{T5}.}
{T5 directly compares encoder--decoder, prefix-LM and decoder-only models under equal budgets: our A2, A3 and A1. It also compares adapters with full fine-tuning (Phases 2--7).}
{Main deck: ``The four architectures'' and ``A controlled experiment plan''.}

\begin{frame}{Roadmap}
  \begin{multicols}{2}\small\tableofcontents\end{multicols}
\end{frame}

% ------------------------------------------------------------------
\section{The problem and the framework}

\begin{frame}{Too many variables, too few controlled comparisons}
  \begin{columns}[T]
    \begin{column}{0.52\textwidth}
      \small
      By 2019, transfer learning (pre-train, then fine-tune) had many competing recipes: BERT, GPT, XLNet, MASS, UniLM\ldots
      \vspace{1mm}

      Each paper changed several things at once: architecture, objective, data, model size, training length. It was unclear \emph{which} choices actually mattered.
      \vspace{1mm}

      \textbf{T5's approach}:
      \begin{enumerate}
        \item One unified \textbf{text-to-text} format for all tasks.
        \item A reasonable baseline.
        \item Change \textbf{one factor at a time} and measure it on many tasks.
        \item Combine the best choices and scale up.
      \end{enumerate}
    \end{column}
    \begin{column}{0.44\textwidth}
      \begin{exampleblock}{Why this matters to us}
        \small
        Our Phase-1 study is the same kind of work on a smaller scale: four architectures, same data, same budget, one variable at a time.
        T5 shows how to design such a study, and gives a prediction for our A1-vs-A2-vs-A3 comparison.
      \end{exampleblock}
    \end{column}
  \end{columns}
\end{frame}

\begin{frame}{The text-to-text framework}
  \begin{columns}[T]
    \begin{column}{0.52\textwidth}
      \small
      Every task becomes: input text (with a short \textbf{task prefix}) $\to$ output text. One model, one loss (maximum likelihood with teacher forcing), one decoding procedure.
      \vspace{1mm}

      \scriptsize
      \begin{tabular}{@{}p{4.3cm}p{1.6cm}@{}}
        \toprule
        Input & Target\\
        \midrule
        translate English to German: That is good. & Das ist gut.\\
        mnli premise: I hate pigeons. hypothesis: My feelings towards pigeons are filled with animosity. & entailment\\
        summarize: \emph{article} & \emph{summary}\\
        stsb sentence1: \ldots\ sentence2: \ldots & 3.8\\
        \bottomrule
      \end{tabular}
    \end{column}
    \begin{column}{0.44\textwidth}
      \begin{exampleblock}{Easy legal examples}
        \small
        \code{judgment: <case text>} $\to$ \code{1}\\
        \code{statutes: <case text>} $\to$ \code{IPC 302, IPC 120B}\\
        \code{role: The accused was arrested on 3 May.} $\to$ \code{Fact}\\
        \code{summarize: <judgment>} $\to$ \code{<summary>}
        \vspace{1mm}

        Even regression (a similarity score such as 3.8) is emitted as text.
      \end{exampleblock}
      \begin{alertblock}{A risk}
        \small
        The model could output an invalid label, e.g.\ ``hamburger''. It counts as wrong; the authors never saw it happen.
      \end{alertblock}
    \end{column}
  \end{columns}
\end{frame}

\begin{frame}{C4: the Colossal Clean Crawled Corpus}
  \begin{columns}[T]
    \begin{column}{0.52\textwidth}
      \small
      Common Crawl produces about 20 TB of text per month, mostly boilerplate, menus, code and gibberish. Filtering heuristics:
      \begin{itemize}
        \item keep lines ending in terminal punctuation; keep lines of 5 or more words; drop pages under 3 sentences;
        \item drop pages with bad words, ``lorem ipsum'', or a curly bracket (code);
        \item drop lines mentioning Javascript or cookie and privacy policies;
        \item deduplicate repeated three-sentence spans;
        \item keep only pages detected as English with probability 0.99 or more.
      \end{itemize}
      The result is about \textbf{750 GB} of clean English text.
    \end{column}
    \begin{column}{0.44\textwidth}
      \begin{exampleblock}{Lessons for our legal corpus}
        \small
        Our Phase-1 data is already legal text, but similar cleaning ideas apply:
        \begin{itemize}
          \item drop OCR fragments and page headers;
          \item deduplicate repeated boilerplate such as cause titles and standard clauses;
          \item drop tiny fragments (our \code{min\_words} filter).
        \end{itemize}
      \end{exampleblock}
    \end{column}
  \end{columns}
\end{frame}

\begin{frame}{The baseline model and objective}
  \begin{columns}[T]
    \begin{column}{0.52\textwidth}
      \small
      \begin{itemize}
        \item \textbf{Encoder--decoder}, each stack like BERT-base: 12 blocks, $d_{\text{model}}=768$, $d_{ff}=3072$, 12 heads. About \textbf{220M} parameters.
        \item Pre-LN with a scale-only LayerNorm (no bias). \textbf{Relative position bias}: a learned scalar per head and per distance bucket, 32 buckets, growing logarithmically up to 128 tokens.
        \item SentencePiece vocabulary of 32k; AdaFactor optimiser; inverse-square-root learning rate.
        \item Pre-training: $2^{19}$ steps $\times$ $2^{16}$ tokens $\approx$ \textbf{34B tokens}. No data is repeated.
      \end{itemize}
    \end{column}
    \begin{column}{0.44\textwidth}
      \begin{exampleblock}{Span-corruption objective}
        \small
        Drop 15\% of tokens. Replace each dropped \emph{span} with a unique sentinel. The target lists only the dropped spans.
        \vspace{1mm}

        Input: ``Thank you \code{<X>} me to your party \code{<Y>} week.''\\
        Target: ``\code{<X>} for inviting \code{<Y>} last \code{<Z>}''
        \vspace{1mm}

        Short targets make pre-training cheap.
      \end{exampleblock}
    \end{column}
  \end{columns}
\end{frame}

\begin{frame}{Does pre-training help, and how noisy are results?}
  \begin{columns}[T]
    \begin{column}{0.56\textwidth}
      \footnotesize
      \begin{tabular}{@{}lcccc@{}}
        \toprule
        & GLUE & CNNDM & SQuAD & SGLUE\\
        \midrule
        Baseline (10 runs) & 83.28 & 19.24 & 80.88 & 71.36\\
        Std.\ deviation & 0.235 & 0.065 & 0.343 & 0.416\\
        No pre-training & 66.22 & 17.60 & 50.31 & 53.04\\
        \bottomrule
      \end{tabular}
      \vspace{2mm}

      \small
      Pre-training gives large gains everywhere except where data is plentiful (e.g.\ EnFr translation, 39.82 vs 39.77).
    \end{column}
    \begin{column}{0.4\textwidth}
      \begin{exampleblock}{A good habit to copy}
        \small
        The authors trained the baseline \textbf{10 times} to measure run-to-run noise. A later change that moves GLUE by only 0.2 is within one standard deviation and may be noise.
        \vspace{1mm}

        For our capstone: repeat the baseline with 2--3 seeds if possible, before claiming a small gain.
      \end{exampleblock}
    \end{column}
  \end{columns}
\end{frame}

% ------------------------------------------------------------------
\section{The systematic study}

\begin{frame}{Architectures compared: it is all about the mask}
  \begin{columns}[T]
    \begin{column}{0.46\textwidth}
      \centering
      \mgrid{fully visible\\(encoder)}{1}\hspace{1mm}
      \mgrid{causal\\(decoder / LM)}{\j<=\i ? 1 : 0}\hspace{1mm}
      \mgrid{causal with prefix\\(prefix LM)}{(\j<3 || \j<=\i) ? 1 : 0}
    \end{column}
    \begin{column}{0.5\textwidth}
      \small
      Structures compared:
      \begin{itemize}
        \item \textbf{Encoder--decoder}: fully visible encoder, causal decoder, cross-attention. Our A2.
        \item \textbf{Encoder--decoder, shared}: the same weights in both stacks.
        \item \textbf{Language model}: one causal stack over input + target. Our A1 and A4.
        \item \textbf{Prefix LM}: one stack; fully visible over the input, causal over the target. Our A3.
      \end{itemize}
      \vspace{1mm}
      \textbf{Fair comparison}: an encoder--decoder with $L+L$ layers has $2P$ parameters but about the \emph{same cost} $M$ as an $L$-layer LM on input + target, because each token passes through only one stack.
    \end{column}
  \end{columns}
\end{frame}

\begin{frame}{Architecture results}
  \centering
  \footnotesize
      \begin{tabular}{@{}llccccc@{}}
        \toprule
        Architecture & Objective & Params & Cost & GLUE & CNNDM & SGLUE\\
        \midrule
        \textbf{Encoder--decoder} & Denoising & $2P$ & $M$ & \textbf{83.28} & \textbf{19.24} & \textbf{71.36}\\
        Enc--dec, shared & Denoising & $P$ & $M$ & 82.81 & 18.78 & 70.73\\
        Enc--dec, 6 layers & Denoising & $P$ & $M/2$ & 80.88 & 18.97 & 68.42\\
        Language model & Denoising & $P$ & $M$ & 74.70 & 17.93 & 55.02\\
        Prefix LM & Denoising & $P$ & $M$ & 81.82 & 18.61 & 68.11\\
        \midrule
        Encoder--decoder & LM & $2P$ & $M$ & 79.56 & 18.59 & 64.29\\
        Language model & LM & $P$ & $M$ & 73.78 & 17.54 & 56.51\\
        Prefix LM & LM & $P$ & $M$ & 79.68 & 17.84 & 64.86\\
        \bottomrule
      \end{tabular}
  \vspace{2mm}
  \begin{exampleblock}{Takeaways}
    \small
    Encoder--decoder wins. Sharing its weights costs little (83.28 to 82.81) while halving parameters.
    \textbf{Prefix LM beats the plain LM} by a wide margin, because the input is fully visible. Denoising beats plain LM objectives.
  \end{exampleblock}
  {\scriptsize Downstream fine-tuning scores, not next-token perplexity; our Phase 1 measures continuation quality, so the ranking may differ.}
\end{frame}

\begin{frame}{Unsupervised objectives}
  \begin{columns}[T]
    \begin{column}{0.56\textwidth}
      \footnotesize
      Variants of BERT-style denoising in the encoder--decoder:
      \vspace{1mm}

      \begin{tabular}{@{}lccc@{}}
        \toprule
        Objective & GLUE & CNNDM & SGLUE\\
        \midrule
        BERT-style & 82.96 & 19.17 & 69.85\\
        MASS-style & 82.32 & 19.16 & 69.28\\
        \textbf{Replace corrupted spans} & 83.28 & 19.24 & \textbf{71.36}\\
        Drop corrupted tokens & \textbf{84.44} & \textbf{19.31} & 68.67\\
        \bottomrule
      \end{tabular}
      \vspace{2mm}

      \small
      Also tested: corruption rates of 10--50\% and average span lengths of 2--10. Differences were small; 50\% corruption was worst.
    \end{column}
    \begin{column}{0.4\textwidth}
      \begin{exampleblock}{Takeaway}
        \small
        Most denoising objectives perform \textbf{similarly}. So choose the one with the \emph{shortest targets}, which trains fastest.
        \vspace{1mm}

        Span corruption with average span length 3 slightly but significantly beat independent token corruption on most non-translation tasks.
      \end{exampleblock}
    \end{column}
  \end{columns}
\end{frame}

\begin{frame}{Pre-training data: domain and size}
  \begin{columns}[T]
    \begin{column}{0.52\textwidth}
      \scriptsize
      \begin{tabular}{@{}lrccc@{}}
        \toprule
        Data set & Size & GLUE & SQuAD & SGLUE\\
        \midrule
        C4 & 745 GB & 83.28 & 80.88 & 71.36\\
        C4, unfiltered & 6.1 TB & 81.46 & 78.78 & 68.04\\
        RealNews-like & 35 GB & 83.83 & 80.39 & 72.38\\
        WebText-like & 17 GB & 84.03 & 81.42 & 71.40\\
        Wikipedia & 16 GB & 81.85 & 81.29 & 68.01\\
        Wikipedia + books & 20 GB & 83.65 & \textbf{82.08} & \textbf{73.24}\\
        \bottomrule
      \end{tabular}
      \vspace{1mm}

      \textbf{Repeating a small data set} over the same 34B tokens:
      \vspace{1mm}

      \begin{tabular}{@{}lccc@{}}
        \toprule
        Unique tokens & Repeats & GLUE & SGLUE\\
        \midrule
        full C4 & 0 & 83.28 & 71.36\\
        $2^{27}$ & 256 & 82.62 & 69.97\\
        $2^{23}$ & 4,096 & 76.34 & 59.29\\
        \bottomrule
      \end{tabular}
    \end{column}
    \begin{column}{0.44\textwidth}
      \begin{exampleblock}{Takeaways}
        \small
        \begin{itemize}
          \item \textbf{Cleaning matters}: unfiltered C4 is worse, although 8$\times$ bigger.
          \item \textbf{In-domain data helps in-domain tasks}. Book text boosted a fiction-based reading task (MultiRC exact match 25.8 to 50.9).
          \item Heavy repetition hurts; the model memorises.
        \end{itemize}
        For us: legal pre-training should help legal tasks. Our budget of about 200M tokens from 4.5B available means \textbf{no repetition} needed.
      \end{exampleblock}
    \end{column}
  \end{columns}
\end{frame}

\begin{frame}{Fine-tuning strategies and multi-task learning}
  \begin{columns}[T]
    \begin{column}{0.52\textwidth}
      \footnotesize
      Updating only part of the model:
      \vspace{1mm}

      \begin{tabular}{@{}lccc@{}}
        \toprule
        Method & GLUE & CNNDM & EnDe\\
        \midrule
        All parameters & \textbf{83.28} & \textbf{19.24} & \textbf{26.98}\\
        Adapters, $d=32$ & 80.52 & 15.08 & 13.84\\
        Adapters, $d=512$ & 81.54 & 17.78 & 23.45\\
        Gradual unfreezing & 82.50 & 18.95 & 26.71\\
        \bottomrule
      \end{tabular}
      \vspace{2mm}

      \small
      \textbf{Multi-task}: training on a mixture of tasks with various mixing rates never quite matched pre-train-then-fine-tune. Multi-task pre-training followed by fine-tuning was comparable.
    \end{column}
    \begin{column}{0.44\textwidth}
      \begin{exampleblock}{For Phases 2--7}
        \small
        \begin{itemize}
          \item Full fine-tuning was best, but most expensive.
          \item Small adapters struggle on \textbf{generation-heavy} tasks (summarisation, translation); enough capacity is needed.
          \item Adapters with $d=32$ are fine for classification (GLUE 80.5) but poor for translation (13.8 BLEU).
        \end{itemize}
        Expect PEFT gaps to be larger in Phase 6 (summarisation) than in Phases 3--5.
      \end{exampleblock}
    \end{column}
  \end{columns}
\end{frame}

\begin{frame}{Scaling: how to spend 4$\times$ more compute}
  \begin{columns}[T]
    \begin{column}{0.54\textwidth}
      \footnotesize
      \begin{tabular}{@{}lccc@{}}
        \toprule
        Strategy & GLUE & SGLUE & EnDe\\
        \midrule
        Baseline & 83.28 & 71.36 & 26.98\\
        1$\times$ size, 4$\times$ steps & 85.33 & 74.72 & 27.08\\
        1$\times$ size, 4$\times$ batch & 84.60 & 74.64 & 27.07\\
        2$\times$ size, 2$\times$ steps & \textbf{86.18} & 77.18 & \textbf{27.52}\\
        4$\times$ size, 1$\times$ steps & 85.91 & \textbf{78.04} & 27.47\\
        4$\times$ ensembled & 84.77 & 71.74 & 28.05\\
        \bottomrule
      \end{tabular}
    \end{column}
    \begin{column}{0.42\textwidth}
      \begin{exampleblock}{Takeaways}
        \small
        \begin{itemize}
          \item Every option helps.
          \item A \textbf{bigger model trained for fewer steps} usually beats a smaller model trained longer.
          \item Ensembles give an independent boost.
        \end{itemize}
        For us: within 50--300M parameters and a fixed Colab budget, test a larger model for fewer steps against a smaller one for more steps.
      \end{exampleblock}
    \end{column}
  \end{columns}
\end{frame}

\begin{frame}{Putting it all together: the T5 models}
  \begin{columns}[T]
    \begin{column}{0.54\textwidth}
      \small
      Final recipe: encoder--decoder, span corruption (average span 3), C4, multi-task pre-training then per-task fine-tuning, and much more compute: over \textbf{1 trillion tokens} in total. Sizes from Small (60M) to 11B.
      \vspace{1mm}

      \footnotesize
      \begin{tabular}{@{}lcc@{}}
        \toprule
        Model & CNN/DM ROUGE-2 & CNN/DM ROUGE-L\\
        \midrule
        Previous best & 20.30 & 40.63\\
        T5-Small & 19.56 & 38.35\\
        T5-Base & 20.34 & 39.40\\
        T5-Large & 20.68 & 39.75\\
        T5-11B & \textbf{21.55} & \textbf{40.69}\\
        \bottomrule
      \end{tabular}
    \end{column}
    \begin{column}{0.42\textwidth}
      \begin{exampleblock}{Headline results}
        \small
        \begin{itemize}
          \item SuperGLUE: \textbf{88.9}, up from 84.6, close to the human score of 89.8.
          \item State of the art on many GLUE, SuperGLUE, SQuAD and CNN/DM measures.
          \item \textbf{Not} state of the art on WMT translation, probably because pre-training was English-only and without back-translation.
        \end{itemize}
      \end{exampleblock}
    \end{column}
  \end{columns}
\end{frame}

% ------------------------------------------------------------------
\section{Assessment and outlook}

\begin{frame}{What T5 predicts for our Phase-1 study}
  \small
  \begin{tabularx}{\textwidth}{@{}lX@{}}
    \toprule
    T5 finding & Hypothesis for our capstone\\
    \midrule
    Encoder--decoder best for fine-tuning & A2 may transfer best to Phases 2--7, even if not best at continuation\\
    Prefix LM $\gg$ plain LM & A3 should beat A1 when the prompt is long and informative\\
    Shared enc--dec nearly as good & Parameter sharing (A3) is a cheap way to get encoder benefits\\
    Denoising $>$ LM objective downstream & Our models all use LM-style objectives; span corruption is a possible extra experiment\\
    In-domain data helps & Legal-only pre-training should help legal tasks; test mixing ratios\\
    Bigger model, fewer steps & Compare 80M and 150M models at equal compute\\
    Run-to-run noise & Repeat key runs with different seeds\\
    \bottomrule
  \end{tabularx}
\end{frame}

\begin{frame}{Strengths and limitations}
  \begin{columns}[T]
    \begin{column}{0.48\textwidth}
      \begin{exampleblock}{Strengths}
        \small
        \begin{itemize}
          \item Very large, carefully controlled study; one variable at a time.
          \item Measures run-to-run variance.
          \item Simple unified format for every task.
          \item Released C4, code and all model checkpoints.
          \item Clear, practical takeaways.
        \end{itemize}
      \end{exampleblock}
    \end{column}
    \begin{column}{0.48\textwidth}
      \begin{alertblock}{Limitations}
        \small
        \begin{itemize}
          \item Enormous compute; hard to reproduce outside large labs.
          \item Factors are studied one at a time, so interactions can be missed.
          \item English-only pre-training.
          \item Judged on fine-tuning benchmarks; open-ended generation quality was not the focus.
          \item Maximum length 512: long documents not studied.
        \end{itemize}
      \end{alertblock}
    \end{column}
  \end{columns}
\end{frame}

\begin{frame}{Future research directions (1): from the paper's outlook and after}
  \small
  \begin{tabularx}{\textwidth}{@{}lX@{}}
    \toprule
    Direction & Details\\
    \midrule
    Stronger small models \emph{(paper)} & Distillation, parameter sharing, conditional computation for low-resource settings\\
    More efficient knowledge extraction \emph{(paper)} & Denoising may be inefficient; ELECTRA (Clark et al., 2020) trains a discriminator instead\\
    Formalise task similarity \emph{(paper)} & A principled way to choose pre-training data for a target task\\
    Language-agnostic models \emph{(paper)} & mT5 (Xue et al., 2021) trains T5 on 101 languages\\
    Instruction tuning & Flan-T5 (Chung et al., 2022) fine-tunes on over 1,800 tasks phrased as instructions\\
    Better objective mixtures & UL2 (Tay et al., 2023) mixes prefix-LM and several span-corruption denoisers\\
    Tokenizer-free T5 & ByT5 (Xue et al., 2022) works directly on bytes\\
    \bottomrule
  \end{tabularx}
\end{frame}

\begin{frame}{Future research directions (2): open questions}
  \begin{columns}[T]
    \begin{column}{0.48\textwidth}
      \begin{block}{Open in general}
        \small
        \begin{itemize}
          \item Which architecture is best when the goal is open-ended generation \emph{and} fine-tuning?
          \item How factors interact: is the best objective different at large scale?
          \item Long-input text-to-text models.
          \item Cheaper ways to reach the same downstream quality.
        \end{itemize}
      \end{block}
    \end{column}
    \begin{column}{0.48\textwidth}
      \begin{exampleblock}{Open for legal NLP}
        \small
        \begin{itemize}
          \item One text-to-text legal model for all seven phases, using task prefixes.
          \item Text outputs for span tasks (NER) need exact character offsets, a known weak spot of text-to-text.
          \item Legal-specific cleaning heuristics, the equivalent of C4 for judgments.
        \end{itemize}
      \end{exampleblock}
    \end{column}
  \end{columns}
\end{frame}

\begin{frame}{Summary}
  \begin{itemize}
    \item \textbf{Text-to-text}: every task is input text $\to$ output text, trained with one loss.
    \item \textbf{C4}: 750 GB of heuristically cleaned web text.
    \item Controlled findings: encoder--decoder best; prefix LM much better than plain LM; denoising beats LM objectives downstream; objectives are similar, so use short targets; clean and in-domain data help; heavy repetition hurts; full fine-tuning beats adapters; bigger models beat longer training.
    \item Scaled to 11B parameters: state of the art on many benchmarks, SuperGLUE 88.9 vs human 89.8.
    \item For us: a template for a controlled architecture study, and concrete predictions for A1--A3.
  \end{itemize}
\end{frame}

\begin{frame}{Reference}
  \small
  Raffel, C., Shazeer, N., Roberts, A., Lee, K., Narang, S., Matena, M., Zhou, Y., Li, W., and Liu, P.\ J.\ (2020). Exploring the Limits of Transfer Learning with a Unified Text-to-Text Transformer. \emph{Journal of Machine Learning Research}, 21(140):1--67. arXiv:1910.10683.
  \vspace{3mm}

  {\footnotesize Later work cited: Clark et al.\ (2020) ELECTRA; Xue et al.\ (2021) mT5; Chung et al.\ (2022) Scaling Instruction-Finetuned Language Models (Flan-T5); Xue et al.\ (2022) ByT5; Tay et al.\ (2023) UL2.}
\end{frame}

\end{document}
